\documentclass[11pt]{amsart}

\usepackage[T1]{fontenc}
\usepackage{lmodern}
\usepackage{microtype}
\usepackage{mathtools,amssymb,amsthm}
\usepackage{array}
\usepackage[hidelinks]{hyperref}

\hypersetup{
  pdftitle={An explicit 32-point set with a free half-turn and no convex 7-gon}
}

\newtheorem{theorem}{Theorem}
\newtheorem{proposition}[theorem]{Proposition}
\newtheorem{corollary}[theorem]{Corollary}
\theoremstyle{remark}
\newtheorem*{remark}{Remark}

\title[A 32-point set with a free half-turn and no convex 7-gon]
{An explicit 32-point set with a free half-turn and no convex 7-gon}

\date{}

\subjclass[2020]{52C10, 52A10, 05D10}
\keywords{Erd\H{o}s--Szekeres problem, convex position, point configuration,
rotational symmetry, explicit construction}

\begin{document}

\begin{abstract}
Subercaseaux, Mackey, Qian and Heule closed the $4$-fold symmetric case of the
$32$-point Erd\H{o}s--Szekeres problem negatively and listed, as a further
problem in their conclusions, the realization of a \emph{$2$-fold} symmetric
$32$-point set with no convex $7$-gon. We exhibit one. The set
$W^{*}=P\cup(-P)$, where $P$ is an explicit list of sixteen integer points of
absolute value at most $31$, lies in general position, is carried to itself by
the half-turn $h(z)=-z$, which acts on it freely with sixteen orbits of size
two, and contains no seven points in convex position; its largest subsets in
convex position have size $6$, and there are $56261$ of them. Realizability is
immediate because the points are integers. Combined with the negative $4$-fold
result and an orbit argument, this determines which rotational orders can act on
a $7$-gon-free set of $32$ points: only $s\in\{1,2,4\}$ are possible, $s=1$ and
$s=2$ occur, and $s=4$ does not. Nothing here bounds $g(7)$.
\end{abstract}

\maketitle

\section{The problem}\label{sec:problem}

Write $g(k)$ for the least $n$ such that every set of $n$ points in general
position in the plane (no three collinear) contains $k$ points in convex
position. Erd\H{o}s and Szekeres proved $g(k)\ge 2^{k-2}+1$ and conjectured
equality \cite{ES35,ES61}; the only known values are $g(4)=5$, $g(5)=9$ and
$g(6)=17$ \cite{SzekeresPeters}, so $g(7)$ is open and $32=2^{7-2}$ is the
conjectured extremal size, not a proved one.

Subercaseaux, Mackey, Qian and Heule \cite{SMQH} settle the case $n=32$, $k=7$
under a $4$-fold rotational symmetry assumption negatively, and their
Conclusions name what is left. Quoted from the e-print:

\begin{quote}
Further problems in this line of work include finding a realization of a
\verb|$2$|-fold symmetric construction of \verb|$32$| points without a convex
\verb|$7$|-gon, and designing an efficient solver that can handle collinearity
constraints for the realizability problem.
\end{quote}

\noindent
The sentence carries no number: it is the further-problems sentence of a
Conclusions section, and it is reproduced verbatim so that what is answered
below does not depend on a label. The ambient definitions are those of the
source. General position is no three points collinear; a set $S$ in general
position is in convex position if and only if every $4$-point subset of $S$ is,
equivalently if and only if no point of $S$ lies in the convex hull of three
others, a consequence of Carath\'eodory's theorem. Here the exhibited object is
a list of integers, so realizability needs no appeal to the existential theory
of the reals.

Constructions of $2^{k-2}$-point $k$-gon-free sets, generalizing the previously
known families, are given by Baek and Balko \cite{BaekBalko}; symmetry is not
considered there. No priority or novelty is claimed for the set below.

\section{The construction}\label{sec:result}

Let $P$ be the sixteen integer points
\[
\begin{array}{cccc}
 (18,-4) & (6,-18) & (-17,4) & (-3,-18)\\
 (-5,20) & (12,21) & (15,23) & (-13,5)\\
 (20,-25) & (-22,2) & (23,27) & (-22,-26)\\
 (23,-2) & (-30,-31) & (27,-1) & (4,19)
\end{array}
\]
and put $W^{*}=P\cup(-P)$, that is, the thirty-two points
\[
\begin{array}{cccc}
 (18,-4) & (-18,4) & (6,-18) & (-6,18)\\
 (-17,4) & (17,-4) & (-3,-18) & (3,18)\\
 (-5,20) & (5,-20) & (12,21) & (-12,-21)\\
 (15,23) & (-15,-23) & (-13,5) & (13,-5)\\
 (20,-25) & (-20,25) & (-22,2) & (22,-2)\\
 (23,27) & (-23,-27) & (-22,-26) & (22,26)\\
 (23,-2) & (-23,2) & (-30,-31) & (30,31)\\
 (27,-1) & (-27,1) & (4,19) & (-4,-19)
\end{array}
\]
Every coordinate is an integer of absolute value at most $31$, the bound being
attained at $(-30,-31)$. That bound is not claimed to be minimal.

\begin{theorem}\label{thm:main}
$W^{*}$ is a set of $32$ points in general position in the plane; the half-turn
$h(z)=-z$ satisfies $h(W^{*})=W^{*}$, has order $2$, and acts on $W^{*}$
freely, with exactly $16$ orbits of size $2$; and $W^{*}$ contains no $7$
points in convex position. Its largest subsets in convex position have size
exactly $6$, and there are $56261$ of them.
\end{theorem}

Thus a realization of the object named in the sentence quoted above exists.

\begin{proof}
The list has $32$ entries and they are pairwise distinct, so $|W^{*}|=32$; the
origin does not occur in the list. The map $h$ is linear, hence an affine
bijection of the plane, so it preserves collinearity, convex hulls and convex
position; $h(W^{*})=W^{*}$ holds entry by entry by construction of $W^{*}$ as
$P\cup(-P)$; $h\circ h=\mathrm{id}$ and $h\neq\mathrm{id}$, so $h$ has order
exactly $2$. The unique fixed point of $h$ in the plane is the origin, which is
not in $W^{*}$; therefore no point of $W^{*}$ is fixed, every orbit has size
$2$, and there are $32/2=16$ of them.

General position and the absence of seven points in convex position are finite
decisions, and no short argument for them is claimed; they were computed as
described in \S\ref{sec:verify}. All $\binom{32}{3}=4960$ orientation
determinants $\det\bigl(b-a,\;c-a\bigr)$ are nonzero, their absolute values
ranging from $1$ to $2892$. For the second part, the subsets of $W^{*}$ in
convex position, counted by size, are
\[
\begin{array}{r|ccccc}
 \text{size} & 3 & 4 & 5 & 6 & \ge 7\\\hline
 \text{number in convex position} & 4960 & 24322 & 59002 & 56261 & 0\\
 \binom{32}{\text{size}} & 4960 & 35960 & 201376 & 906192 &
\end{array}
\]
so the maximum size of a subset in convex position is $6$ and the number of
those is $56261$. Convex position is inherited by subsets, so ``no $7$ in
convex position'' already gives ``no $m$ in convex position'' for every
$m\ge 7$.
\end{proof}

\begin{remark}
Two entries of the table are forced: every triple of a set in general position
is in convex position, which fixes the size-$3$ column, and $W^{*}$ must
contain a convex hexagon since $g(6)=17\le 32$. What is not forced is the $0$
in the final column, that is, that the maximum is \emph{exactly} $6$. An
explicit hexagon, for a reader who wants one to hand, is
\[
 (18,-4)\quad(-18,4)\quad(6,-18)\quad(-6,18)\quad(-3,-18)\quad(3,18).
\]
\end{remark}

\section{Rotation orders at $n=32$}\label{sec:complete}

\begin{proposition}\label{prop:orders}
Let $W$ be a set of $32$ points in general position with no $7$ points in
convex position, and let $\rho$ be a rotation of the plane of order $s$ with
$\rho(W)=W$. Then $s\in\{1,2,4\}$.
\end{proposition}

\begin{proof}
The cyclic group $\langle\rho\rangle$ has order $s$ and acts on $W$; the only
point of the plane with nontrivial stabiliser is the centre of $\rho$, so every
orbit has size $s$ except for at most one orbit of size $1$. Hence
$32\equiv 0$ or $1 \pmod s$, i.e.\ $s$ divides $32$ or $s$ divides $31$, which
leaves $s\in\{1,2,4,8,16,31,32\}$. In particular $s\in\{3,5,6\}$ is impossible,
since $32\equiv 2$ modulo each of $3$, $5$ and $6$. Finally, if $s\ge 7$ then
$W$ contains an orbit of size $s$, which is the vertex set of a regular
$s$-gon; any $7$ of its vertices are in convex position, a contradiction. This
deletes $8,16,31,32$ and leaves $\{1,2,4\}$.
\end{proof}

\begin{corollary}\label{cor:complete}
For $n=32$ and $k=7$: a set with $s=1$ exists (Erd\H{o}s--Szekeres
\cite{ES61}); a set with $s=2$ exists (Theorem~\ref{thm:main}); no set with
$s=4$ exists (\cite{SMQH}, established there by an exhaustive computation we do
not reproduce); and no other rotational order is possible
(Proposition~\ref{prop:orders}).
\end{corollary}

\noindent
Of the four lines, only the $s=2$ line is new here; the $s=4$ line is the
theorem of \cite{SMQH} and the $s=1$ line is classical. The corollary concerns
rotations only: the mirror-symmetric case at $n=32$ is untouched. Nor does
anything above bound $g(7)$, which is unknown; the inequality $g(7)\ge 33$ is
classical, from the asymmetric $2^{k-2}$ construction of \cite{ES61} (see also
\cite[Thm.~2.6]{MorrisSoltan}), and $W^{*}$ does not improve it.

\section{Verification}\label{sec:verify}

The finite parts of Theorem~\ref{thm:main} were recomputed by two structurally
different exhaustive procedures, in exact integer arithmetic with no
floating-point decision anywhere. The first applies the Carath\'eodory
criterion of \S\ref{sec:problem} to each of the $\binom{32}{7}=3{,}365{,}856$
seven-subsets, with no pruning, and finds $0$ in convex position. The second
never mentions the number $7$: it enumerates every subset in convex position
\emph{at every size at once}, by growing angularly ordered left-turning chains
from each vertex in turn with the closing turn checked explicitly, and so
returns the maximum directly; because that enumeration is untruncated, its
answer $6$ settles all sizes above $6$ simultaneously. The two agree on the
size-$4$ column, where both are cheap: $24322$ of $\binom{32}{4}=35960$. The
criterion was pinned in both polarities before use, so that neither answer is
vacuous.

A program that carries out all of this on the integers printed in
\S\ref{sec:result}, together with a recorded run of it, accompanies this note;
it is not needed in order to redo the work, since every input it consumes is
printed above. It does \emph{not} re-run the search that produced $W^{*}$,
which is why $W^{*}$ is pinned as a printed list of integers, and it does
\emph{not} recompute the $4$-fold negative result quoted in
Corollary~\ref{cor:complete}, which is taken from \cite{SMQH}.

\begin{thebibliography}{10}

\bibitem{ES35}
P.~Erd\H{o}s and G.~Szekeres,
\emph{A combinatorial problem in geometry},
Compositio Math. \textbf{2} (1935), 463--470.

\bibitem{ES61}
P.~Erd\H{o}s and G.~Szekeres,
\emph{On some extremum problems in elementary geometry},
Ann. Univ. Sci. Budapest. E\"otv\"os Sect. Math. \textbf{3--4} (1960/61),
53--62.

\bibitem{MorrisSoltan}
W.~Morris and V.~Soltan,
\emph{The Erd\H{o}s--Szekeres problem on points in convex position ---
a survey},
Bull. Amer. Math. Soc. \textbf{37} (2000), no.~4, 437--458.
\href{https://doi.org/10.1090/S0273-0979-00-00877-6}%
{doi:10.1090/S0273-0979-00-00877-6}.

\bibitem{SzekeresPeters}
G.~Szekeres and L.~Peters,
\emph{Computer solution to the 17-point Erd\H{o}s--Szekeres problem},
ANZIAM J. \textbf{48} (2006), no.~2, 151--164.
\href{https://doi.org/10.1017/S144618110000300X}%
{doi:10.1017/S144618110000300X}.

\bibitem{SMQH}
B.~Subercaseaux, E.~Mackey, L.~Qian, and M.~J.~H. Heule,
\emph{Automated symmetric constructions in discrete geometry},
in: Intelligent Computer Mathematics (CICM 2025), Lecture Notes in Computer
Science \textbf{16136}, Springer, Cham, 2026, pp.~29--47.
\href{https://doi.org/10.1007/978-3-032-07021-0_3}%
{doi:10.1007/978-3-032-07021-0\_3};
\href{https://arxiv.org/abs/2506.00224}{arXiv:2506.00224}.
The sentence quoted in \S\ref{sec:problem} is from the e-print of v1
(30 May 2025).

\bibitem{BaekBalko}
J.~Baek and M.~Balko,
\emph{The Erd\H{o}s--Szekeres conjecture revisited},
in: 41st International Symposium on Computational Geometry (SoCG 2025),
LIPIcs \textbf{332}, 13:1--13:15.
\href{https://doi.org/10.4230/LIPIcs.SoCG.2025.13}%
{doi:10.4230/LIPIcs.SoCG.2025.13};
full version \href{https://doi.org/10.1016/j.jcta.2026.106195}%
{doi:10.1016/j.jcta.2026.106195}.

\end{thebibliography}

\end{document}
